\section*{Chapitre \ref{ch:s2:cross-currency-basis-and-fx-carry} --- Base de devises et carry de change}

\begin{solution}{exo:s2:cross-currency-basis-and-fx-carry:1}
$2.7/7.2 = 0.375$ avec les chiffres arrondis, 0.37 sans arrondi.
\end{solution}

\begin{solution}{exo:s2:cross-currency-basis-and-fx-carry:2}
$(50 - 15)/252 \times 10 = 1.39$ points de base du montant prêté.
\end{solution}

\begin{solution}{exo:s2:cross-currency-basis-and-fx-carry:3}
$(1.14 - 0.63)/1.30 = 0.39$ avec les chiffres arrondis, 40~\% sans arrondi.
\end{solution}

\begin{solution}{exo:s2:cross-currency-basis-and-fx-carry:4}
Son taux d'intérêt était le plus bas, ou proche du plus bas, des neuf devises sur toute la période, souvent nul ou négatif, si bien qu'il
figurait toujours parmi les trois devises de financement.
\end{solution}

\begin{solution}{exo:s2:cross-currency-basis-and-fx-carry:5}
Elle baissait un peu chaque mois~; un put dont le strike se situe à 1.5 écarts types sous zéro ne paie que si le mois baisse davantage que cela, si bien que les primes
ont été dépensées et peu est revenu.
\end{solution}

\begin{solution}{exo:s2:cross-currency-basis-and-fx-carry:6}
Refermer la position exige du bilan~: l'arbitragiste doit emprunter des dollars directement et les prêter via le swap, une position brute que les banques font payer,
surtout aux dates de reporting. La base se fixe là où ce coût rend la transaction non rentable pour la banque marginale.
\end{solution}

\begin{solution}{exo:s2:cross-currency-basis-and-fx-carry:7}
La couverture coûte désormais 0.55~\% par an et rapporte 0.63~\%~: elle dégage un gain, et le ratio de Sharpe du panier couvert passe à 0.26 sans skew. La prime de krach
n'est pas dans les seuls rendements des devises~; elle se trouve dans ce que les vendeurs d'options facturent au-dessus de la fair value, soit ici 1.2 fois la volatilité passée
retenue par la valorisation du chapitre.
\end{solution}

\begin{solution}{exo:s2:cross-currency-basis-and-fx-carry:8}
Un levier de 2.8 multiplie aussi les krachs~: la perte de l'automne 2008, de 19.4~\%, passe à environ 54~\%. Le ratio de Sharpe masque un skew négatif de $-0.60$~;
dimensionner sur le krach, et non sur la volatilité.
\end{solution}

\begin{solution}{pb:s2:cross-currency-basis-and-fx-carry:1}
\begin{enumerate}
\item Position longue sur les devises à taux les plus élevés, courte sur les plus basses~; à partir des taux de change FRED et des taux interbancaires de l'OCDE, en données mensuelles.
\item 2.7~\% par an pour une volatilité de 7.2~\%, ratio de Sharpe de 0.37, skew de $-0.60$~; $-19.4~\%$ à l'automne 2008, $-10.5~\%$ en octobre 2008.
\item Position longue surtout sur les dollars néo-zélandais et australiens ainsi que la krone~; courte en permanence sur le franc suisse, souvent sur le yen, la krona et l'euro.
\item Le carry a rapporté des ratios de Sharpe comparables à ceux des facteurs actions~; les primes de krach ne représentaient au plus qu'un tiers du rendement.
\item Carry avec des options qui paient en cas de krach.
\item Sans couverture~: 1.30~\% par an, ratio de Sharpe 0.23, skew $-0.39$~; avec couverture~: 0.78~\%, 0.15, 0.01.
\item Les mois des premier et troisième krachs sont ramenés de $-12.0~\%$ à $-7.9~\%$ et de $-14.4~\%$ à $-8.5~\%$~; le deuxième a à peine changé.
\item C'est le prix à payer pour retirer le krach du rendement.
\item Emprunter des dollars directement ou via un swap devrait coûter la même chose~; la base est l'écart.
\item Des écarts importants et persistants, plus forts aux fins de trimestre, qui renvoient à la réglementation.
\item Prêter des dollars via le swap~; 11.2 points de base par an après un coût de 15 points de base.
\item 42.5 points de base par an pendant la durée du prêt, sur 40 jours par an.
\item \textbf{Résultat à retenir.} Le ratio de Sharpe du panier de carry synthétique passe de 0.23 à 0.15 avec des couvertures mensuelles par put, qui retirent son skew négatif
à hauteur de 40~\% de son rendement~; le basis trade rapporte 42.5 points de base par an pendant les prêts sur les fenêtres de fin de trimestre, 6.7 par an en tout.
\item Les points de terme incluent la base, et les krachs du carry sont des événements de financement.
\item Moins de carry dans les périodes calmes, une certaine protection dans les krachs longs, aucune dans les krachs rapides.
\item Par le krach, non par la volatilité.
\item Les cours des swaps pour les dates exactes, et le \omterm{def:s2:swap-spread-and-asset-swap-trades:bscost}{coût de bilan} propre au prêteur.
\item L'arbitrage de funding par la base.
\item Les deux sont des prix du bilan bancaire~: le spread de swap pour la détention d'obligations, la base pour le prêt de dollars.
\item Le coût du bilan bancaire pour prêter des dollars.
\end{enumerate}
\end{solution}

\begin{solution}{iq:s2:cross-currency-basis-and-fx-carry:1}
Emprunter dans les devises à faible taux et prêter dans celles à taux élevé~; en moyenne, les taux de change ne bougent pas assez pour compenser le différentiel, mais ils le font lors des
krachs. Le rendement est la rémunération de la prise de ce risque.
\end{solution}

\begin{solution}{iq:s2:cross-currency-basis-and-fx-carry:2}
Couvrir les positions de carry avec des options hors de la monnaie aux prix de marché et comparer les rendements couverts et non couverts~; l'écart est la prime payée
pour retirer le krach.
\end{solution}

\begin{solution}{iq:s2:cross-currency-basis-and-fx-carry:3}
Se convertir en dollars via un swap coûte 40 points de base par an de plus que la parité~; les emprunteurs de dollars passant par des swaps (banques non américaines, investisseurs couvrant
des actifs libellés en dollars) le paient à celui qui peut prêter des dollars via le swap et dispose du bilan.
\end{solution}

\begin{solution}{iq:s2:cross-currency-basis-and-fx-carry:4}
Une fuite vers les actifs sûrs à la manière de 2008 (devises à haut rendement en baisse de 15--20~\%, devises de financement en hausse), un choc de taux sur une devise de financement et un événement de liquidité qui
élargit les spreads à terme.
\end{solution}

\begin{solution}{iq:s2:cross-currency-basis-and-fx-carry:5}
Cours spot et points de terme par échéance, taux de dépôt, calendriers des jours fériés par devise, dates de valeur~; la fin de mois apporte les dates de roll, les jours fériés qui décalent les
dates de valeur et les cotations peu liquides.
\end{solution}

\begin{solution}{iq:s2:cross-currency-basis-and-fx-carry:6}
Emprunter un dollar au taux $r_\$$ ou le convertir au cours spot $S$ (devise étrangère par dollar), prêter au taux $r_f$ et reconvertir au cours à terme $F$ doit donner le même résultat~:
$F = S(1 + r_f\tau)/(1 + r_\$\tau)$. Avec une base, $F = S(1 + (r_f + b)\tau)/(1 + r_\$\tau)$~: un $b$ négatif rend la jambe étrangère moins valorisée, si bien que les dollars
obtenus via le swap coûtent plus cher.
\end{solution}
